\section{Split-invisible classes and the global firewall}
\label{sec:counterpacket}

Fix a coefficient prime \(\ell\), a finite extension
\(E_\ell/\Ql\) containing the traces in scope, and a finite set \(S\)
containing \(\ell\), the ramified primes, and the bad primes of every
compatible system under discussion.  For \(F=\Q\) or \(K\), let \(S_F\)
be the primes of \(F\) above \(S\), and let
\[
K_{0,\ell}^{\semis}(F;S_F)
\]
be the Grothendieck group generated by the finite-dimensional continuous
semisimple \(E_\ell\)-representations of \(G_F\) arising from these
fixed-\(\ell\) compatible-system realizations and their semisimple
subquotients, all unramified outside the primes above \(S\).  This is not
the unrestricted category of all representations of \(G_F\).  Restriction
is the map
\[
\Res:K_{0,\ell}^{\semis}(\Q;S)
\longrightarrow K_{0,\ell}^{\semis}(K;S_K).
\]

\begin{proposition}[Counterpacket firewall]
\label{prop:counterpacket}
Let \(D\) be a virtual rational class in the preceding fixed-\(\ell\),
finite-ramification category.  If its trace vanishes at all but a
relative-Dirichlet-density-zero subset of the good rational primes split in
\(K\), then
\[
\Res(D)=0.
\]
If \(D\) is represented by an actual system, then \(D=0\).
\end{proposition}

\begin{proof}
Prime ideals of \(K\) of residue degree one form a
Dirichlet-density-one set.  Outside the finite ramified set they lie over
split rational primes, and the exceptional rational split primes in the
hypothesis lift to a density-zero subset of that set.  Hence the trace of
\(\Res(D)\) vanishes on a dense set of Frobenius conjugacy classes.
Chebotarev and Brauer--Nesbitt imply \(\Res(D)=0\).  If \(D\) is
actual, restriction preserves its nonnegative rank; a zero restricted class
therefore has rank zero and is the zero representation.
\end{proof}

Restriction is not injective on virtual classes.  The explicit class
\begin{equation}
\one-\chi_{K/\Q}
\label{eq:explicit-kernel-class}
\end{equation}
is nonzero over \(\Q\) because \(K/\Q\) is a nontrivial quadratic
extension, but it restricts to zero over \(K\).  More generally,
\[
U-U\otimes\chi_{K/\Q}
\]
is a kernel class, and it is nonzero when
\(U\not\simeq U\otimes\chi_{K/\Q}\).  Every such kernel class has rank
zero.  It can change a rational extension and its inert traces, but it
cannot change a \(K\)-rail rank or a common \(G_3\)-isotypic
multiplicity.  Thus virtual split-invisible counterpackets do not evade
Theorem~\ref{thm:denominator} or
Corollary~\ref{cor:no-central-sector}.

\subsection{Why split-local identities do not globalize}

At a good inert rational prime \(p\), let \(v\) be the prime of \(K\)
above \(p\), and write
\[
P_p(U)=\prod_i(1-\alpha_iU).
\]
Because \(F_v=F_p^2\), one has the exact identity
\begin{equation}
P_{K,v}(U^2)
=\prod_i(1-\alpha_i^2U^2)
=P_p(U)P_p(-U).
\label{eq:inert-polynomial}
\end{equation}
The right-hand side is generally not a square.  The equality of the two
degree-one factors at a split rational prime therefore supplies no
corresponding inert square root.

The conclusions of this paper are accordingly confined to good split
primes.  They do not construct a global or inert fractional Euler root,
prove smoothness or packet-admissibility for \(n\ge5\), supply bad-place
data, or establish automorphy, meromorphic continuation, a functional
equation, or RH.
